\chapter*{模板使用说明}
\addcontentsline{toc}{chapter}{模板使用说明}
本模板根据本目录的 \path{problem.md} 及相邻 \path{../Pgl50h/} 工程整理，
适用于“基于紫光同创 FPGA 的任意角度的畸变矫正”选题。
模板以 2026 年 10 月 22 日 CP2102 位流及其归档证据为主要基线，
预填系统架构、算法、模块职责和已有验证数据。所有蓝色“待补充”字段、撰写提示和图片框，
均应在正式提交前完成替换或删除。

本文使用三种结果口径：\textbf{工程实现}描述源码中已有的功能；
\textbf{仿真/实现工具结果}引用已有 RTL 日志及 PDS 报告；
\textbf{实板测量}留待实际硬件验证。当前归档身份文件明确标记
\path{board_programmed=false}、\path{board_tested=false}。
模板整理过程中仅查阅工程与验证记录，未重新运行 FPGA 综合、布局布线或 RTL 回归。

报告中的“任意角度”沿用赛题名称；当前基线仅使用 yaw/pitch，默认每轴约
$\pm15^\circ$，roll 不参与变换。默认实时分块路径使用最近邻采样；
历史串行路径保留双线性，但其性能不适用本模板中的分块吞吐数据。
1080p60 表示输入/输出扫描时序，1080p30 表示当前行为模型仿真的新矫正帧更新能力。

\guide{首先填写封面和第 1 章指标，再补充第 6、7 章实测数据。
如更换源码、标定或位流，应同时更新基线、结果表和附录证据，不能沿用旧性能结论。}

\chapter*{摘要}
\addcontentsline{toc}{chapter}{摘要}
针对镜头畸变和相机姿态变化引起的图像几何失真，本文基于紫光同创
PGL50H-FBG484-6 FPGA，设计面向 $1920\times1080$ 视频的几何矫正系统。
系统通过 MS7200 接收 HDMI RGB888 数据，在量化为 RGB565 前完成帧自适应
Gamma 预处理，并将源帧写入 DDR3。矫正器根据相机标定内参、二阶径向畸变系数及
逐帧锁存的 IMU yaw/pitch，通过定点逆投影生成采样坐标，采用
$128\times32$ 双槽分块缓存和突发访问生成结果帧，经 MS7210 输出 HDMI 视频。
板载 CP2102 提供去畸变、姿态补偿和 Gamma 的独立控制及亮度统计遥测。

当前基线在包含并发采集、渲染、显示及 DDR 停顿的完整 RTL 行为模型中，
达到约 30.000030 个新结果帧/秒，最慢观测渲染时间为 33.14614 ms；
显示扫描为 60 Hz。PDS 2025.2 最终双工作角时序检查通过，
最终占用 104/134 个 DRM 位置、16,878/42,800 个 LUT（含常量和旁路）。
上述数据来自归档仿真及实现报告；实板端到端延迟、真实帧率、直线度及
MTF50/SFR 等结果为\fillin{实测结果及最终结论}。

\noindent\textbf{关键词：}紫光同创 FPGA；畸变矫正；逆投影；IMU；分块缓存；自适应 Gamma

\guide{正式摘要建议保留“问题—方法—结果—结论”结构。
补齐实测后更新结果段，并保持与第 7 章一致；不能把渲染耗时写成端到端延迟。}
